\section{Reproducibility Record}
\label{app:repro}

Our retained experiment record contains the exact harness JSON for every
incumbent and candidate, immutable instance manifests, per-attempt JSONL records, execution
configuration, verifier output, trace paths, and SHA-256 trace digests. Runs use
accounting schema \texttt{shrlm-accounting/v2}. Each record stores instance ID,
attempt number, pass/fail cause, primary quality, input and output tokens,
provider cost, execution time, and whether usage is a lower bound. Tables in the
paper are computed from these records without dropping terminal failures. An
anonymized release artifact should preserve this structure and remove provider
credentials, local paths, and author-identifying metadata.

\paragraph{Frozen-test aggregation.}
For each method and task, we average its two attempts; the reported overall F1
is the equally weighted mean over tasks. Length-specific means use the 16 tasks
in that stratum. Total tokens and cost sum all attempts. A terminal failure is
an attempt whose cause is a resource termination or unparsable final format.
Paired bootstrap replicates sample task IDs with replacement and retain both
methods' attempts for each sampled task. Percentile endpoints use 200,000
replicates with a fixed analysis seed.

\paragraph{Artifact identities.}
The compared generic harnesses are the sparse round-zero harness, the pinned
upstream-prompt reference, and the promoted incumbent saved after round 10.
The specialized reference is the pinned $\lambda$-RLM reconstruction used in
both frozen OOLONG tests. The 48-task manifest is identical across all four
methods and contains 16 tasks at each nominal context length. The 262K manifest
contains 40 tasks, and the OBLIQ manifest contains 40 queries; each manifest is
reused verbatim across compared methods.

\paragraph{Decision semantics.}
The optimization objective is macro F1. Per-length non-regression and an
increase in \texttt{wrong\_format} are vetoes. Recorded token and provider-cost
changes are audit metrics, not a scalarized objective, and the configured cost
band is not enforced for this run. Rounds 1--5 use improvement threshold zero;
rounds 6--10 use .01, with round six recomputed after the change. No earlier
promotion is rewritten. These facts are necessary to reproduce the incumbent
sequence in Figure~\ref{fig:optimization-trajectory}.

\paragraph{What the artifact does not establish.}
OpenRouter does not guarantee replay-identical stochastic completions. The
artifact therefore supports verification of inputs, code paths, accounting,
and reported aggregation, but an independent rerun should be treated as a
replication rather than a bitwise reproduction. Provider charges exclude local
CPU work and should not be interpreted as complete energy or systems cost.
